\begin{frame}
    \frametitle{Положения, выносимые на защиту}
    \begin{itemize}
        \item Методы построения признакового описания для древнекитайского иероглифа в виде конечного набора критичных точек на основе медиального представления и непрерывной морфологии. Предложен метод выбора критических точек среди вершин скелетного графа и точек максимальной кривизной ребер скелетного графа.
        \item Методы построения оптимизационной модели для сравнения точечных признаков иероглифа. Разработаны методы получения меры близости с двумя видами моделей путем решения задачи о назначениях, которые могут быть реализованы эффективной процедурой вычисления.
        \item Обоснование работоспособности предложенных методов путём реализации программного комплекса. Экспериментально доказано, что разработанные методы не уступают по качеству существующим современным методам с использованием нейронных сетей, при этом обладают преимуществом отсутствия требования обучения на гигантской обучающей выборке, низкого запроса к вычислительным ресурсам и интерпретируемости алгоритма.
    \end{itemize}
\end{frame}
\note{
    Проговариваются вслух положения, выносимые на защиту
}

\begin{frame}
    \frametitle{Научная новизна}
    \begin{itemize}
        \item Предложены новые математические модели, позволяющие эффективно анализировать форму древного иероглифа на основе непрерывной морфологии, и построить признаковое описание для изображения иероглифа в виде конечного множества критичных точек, характеризующих топологическое и геометрическое свойство иероглифа.
        \item Предложены модели для сравнения иероглифических признаковых описаний, полученных предложенными моделями построения точечных признаков, по принципу сопоставления соответствия пар критичных точек.
        \item Разработана оптимизационная модель через задачу линейного программирования для вычисления количественной меры сходства между иероглифами на основе точечных признаковых описаний.
    \end{itemize}
\end{frame}
\note{
    Проговаривается вслух научная новизна
}

\begin{frame}
    \frametitle{Научная и практическая значимость}
    \begin{itemize}
        \item Разработаны методы построения признакового описания древнекитайского рукописного иероглифа в виде конечного набора критичных точек через медиальное представление формы бинарного изображения.
        \item Разработаны методы сравнения форм на основе оптимизационной модели линейного программирования, используя построенные точечные признаки иероглифов.
        \item Реализованы все предложенные методы, а также исследованы приложения для решения различных прикладных задач, связанных с поиском информации в архивных документах. Подтверждена сравнимость результатов с передовыми нейросетевыми методами при минимальным требование к ресурсам.
    \end{itemize}
\end{frame}
\note{
    Проговариваются вслух научная и практическая значимость
}

% \begin{frame}
%     \frametitle{Свидетельство о регистрации программы}
%     \begin{figure}[h]
%         \centering
%         \includegraphics[height=0.7\textheight]{registration}
%     \end{figure}
% \end{frame}
% \note{
%     Получено свидетельство о регистрации разработанной программы \textsc{Hello~world™}.
% }

% \begin{frame}
%     \frametitle{Акт о внедрении}
%     \begin{figure}[h]
%         \centering
%         \fbox{
%             \begin{minipage}[t]{0.4\linewidth}
%                 \includegraphics[width=\linewidth]{implementation}
%             \end{minipage}
%         }
%     \end{figure}
% \end{frame}
% \note{
%     Получен акт о внедрении.
% }

\begin{frame} % публикации на одной странице
    % \begin{frame}[t,allowframebreaks] % публикации на нескольких страницах
    \frametitle{Основные публикации}
    \begin{itemize}
        \item  M. Xu, L. Mestetskiy. An Approach to Searching for Hieroglyphs in Chinese Manuscript Archives Based on Morphological Analysis // The International Archives of the Photogrammetry, Remote Sensing and Spatial Information Sciences. — 2025 \textbf{(Scopus)}. — Т. XLVIII-2/W9—2025. — С. 283—288. — URL: https://isprs-archives.copernicus.org/articles/XLVIII-2-W9-2025/283/2025/.
        \item L. Mestetskiy, M. Xu. Constellation method for one-shot search of handwrittenglyphs // Pattern Recognition and Image Analysis. — 2026 \textbf{(В печати на 2-й номер, ВАК, Scopus, WoS)}.
        \item M. Xu, L. Mestetskiy. Handwritten Chinese character retrieval via constellation matching  // Proceedings of the 3rd International Conference on Machine Intelligence and Digital Applications. — 2026 \textbf{(В печати (к мае), Scopus)}.
    \end{itemize}
\end{frame}

\begin{frame} % публикации на одной странице
    % \begin{frame}[t,allowframebreaks] % публикации на нескольких страницах
    \frametitle{Основные публикации}
    \begin{itemize}
        \item М. Сюй. Применение морфологического анализа для поиска и распознавания в рукописных древнекитайских архивах  // Сборник тезисов Международной научной конференции студентов, аспирантов и молодых учёных «Ломоносов 2025» секция «Вычислительная Математика и Кибернетика ». — Макс Пресс. 2025. — С. 74—75.
        \item Л. Местецкий, М. Сюй. Поиск рукописных китайских иероглифов методом созвездий // Тезисы докладов 22-й всероссийской конференции с международным участием «Математические методы распознавания образов». — Макс Пресс. 2025. — С. 112—115.
    \end{itemize}
\end{frame}

\begin{frame}
    \frametitle{Участие в конференциях}
    \begin{itemize}
        \item Международной научной конференции студентов, аспирантов и молодых учёных <<Ломоносов-2025>> (Россия, Москва, 2025),
        \item 6th International workshop on Photogrammetric techniques for environmental and infrastructure monitoring, Biometry and Biomedicine (Россия, Москва, 2025),
        \item 22-й всероссийской конференции с международным участием «Математические методы распознавания образов» (ММРО-2025),
        \item 3rd International Conference on Machine Intelligence and Digital Applications (Китай, Сиань, 2026).
    \end{itemize}
\end{frame}
\note{
    Работа была представлена на ряде конференций.
}

\begin{frame}[plain, noframenumbering] % последний слайд без оформления
    \begin{center}
        \Huge
        Спасибо за внимание!
    \end{center}
\end{frame}
